\documentclass[11pt]{article}
\usepackage[margin=1in]{geometry}
\usepackage{amsmath,amssymb,amsthm}
\usepackage{booktabs}
\usepackage{hyperref}

\newtheorem{theorem}{Theorem}
\newtheorem{proposition}{Proposition}
\newtheorem{remark}{Remark}

\newcommand{\Q}{\mathbb{Q}}
\newcommand{\F}{\mathbb{F}}

\title{A CRT-Controlled 17/19 Point-Count Exchange\\
in a Rank-Positive Elliptic Curve Family}
\author{BCS Research}
\date{\today}

\begin{document}
\maketitle

\begin{abstract}
We study the one-parameter family of elliptic curves
\[
E_t:\quad y^2=x^3+(17+t)x^2+5x+4.
\]
We prove an exact local congruence phenomenon: for every integer
\(t\equiv 66 \pmod{323}\), the reductions satisfy
\[
\#E_t(\F_{17})=19,\qquad \#E_t(\F_{19})=17.
\]
The family also has a non-torsion rational section \(P=(0,2)\), so its
generic Mordell--Weil rank over \(\Q(t)\) is at least one. Sage computations
inside the same CRT class find certified rank-three specializations at
\(t=66\) and \(t=1681\), both with matching analytic rank three.
\end{abstract}

\section{The family}

Let
\[
E_t:\quad y^2=x^3+(17+t)x^2+5x+4.
\]
The point \(P=(0,2)\) lies on every fiber, since substituting \(x=0\)
gives \(y^2=4\).

\section{The CRT point-count exchange}

\begin{theorem}
For every integer \(t\equiv 66 \pmod{323}\),
\[
\#E_t(\F_{17})=19,\qquad \#E_t(\F_{19})=17.
\]
\end{theorem}

\begin{proof}
Over \(\F_p\), the equation only depends on \(t\bmod p\). A direct residue
count gives
\[
t\equiv 15\equiv -2 \pmod{17}\quad\Longrightarrow\quad
\#E_t(\F_{17})=19,
\]
and
\[
t\equiv 9 \pmod{19}\quad\Longrightarrow\quad
\#E_t(\F_{19})=17.
\]
Solving
\[
t\equiv -2 \pmod{17},\qquad t\equiv 9 \pmod{19}
\]
by the Chinese Remainder Theorem gives \(t\equiv 66\pmod{323}\).
\end{proof}

\section{Generic rank}

\begin{proposition}
The generic Mordell--Weil rank of \(E/\Q(t)\) is at least one.
\end{proposition}

\begin{proof}
The point \(P=(0,2)\) is a rational section. If it were torsion over
\(\Q(t)\), then its specialization at every smooth fiber would be torsion.
At \(t=66\), Sage verifies that \(E_{66}(\Q)\) has trivial torsion and rank
three; in particular, \(P_{66}=(0,2)\) is non-torsion. Thus \(P\) is
non-torsion over \(\Q(t)\).
\end{proof}

\section{Rank evidence}

For tested values \(t=66+323k\), Sage gives the following data.

\begin{center}
\begin{tabular}{rrrr}
\toprule
\(t\) & rank bounds & rank & analytic rank\\
\midrule
-1226 & (1,1) & 1 & 1\\
-903  & (1,1) & 1 & 1\\
-580  & (2,2) & 2 & 2\\
-257  & (1,1) & 1 & 1\\
66    & (3,3) & 3 & 3\\
389   & (2,2) & 2 & 2\\
712   & (2,2) & 2 & 2\\
1035  & (1,1) & 1 & 1\\
1358  & (2,2) & 2 & 2\\
1681  & (3,3) & 3 & 3\\
\bottomrule
\end{tabular}
\end{center}

For \(t=66\), Sage finds generators
\[
(-1,9),\quad (0,2),\quad (-45,277).
\]
For \(t=1681\), Sage finds generators
\[
(0,2),\quad (153/4,12751/8),\quad (12/121,6128/1331).
\]

\section{Caveats}

The CRT point-count exchange is exact. The rank pattern is computational
evidence from tested specializations; we do not claim that all
\(t\equiv 66\pmod{323}\) have positive rank, nor that the rank is bounded
below by two or three on the entire CRT class.

\begin{remark}
This note is mathematical pedigree for the BCS research program. It is not a
cryptographic security proof for BCS-521, whose security argument concerns a
separate 521-bit prime-field elliptic-curve group.
\end{remark}

\end{document}

